\documentclass[twocolumn]{aastex631}
% Auto-generated by jansky_research.vgpra._write_macros -- do not edit.
% vgSyn* = synthetic recover-a-known (always live); vgReal{U,N}* = real Uranus/Neptune legs.
\newcommand{\vgSource}{PDS-PPI Voyager 2 PRA low-band 6-s encounter volumes (Uranus + Neptune)}
\newcommand{\vgSynInjected}{17.24}
\newcommand{\vgSynRecovered}{17.22252}
\newcommand{\vgSynSigma}{0.02049}
\newcommand{\vgSynRayPeriod}{17.24308}
\newcommand{\vgSynRecoveredOK}{yes}
\newcommand{\vgSynRayOK}{yes}
\newcommand{\vgRealUPeriod}{18.44362}
\newcommand{\vgRealUSigma}{1.94272}
\newcommand{\vgRealULo}{14.41862}
\newcommand{\vgRealUHi}{18.67411}
\newcommand{\vgRealURayPeriod}{14.72424}
\newcommand{\vgRealUSpan}{312.0}
\newcommand{\vgRealUUnc}{1.94272}
\newcommand{\vgRealUSpread}{1.75737}
\newcommand{\vgRealURecovers}{no}
\newcommand{\vgRealUHist}{17.24}
\newcommand{\vgRealUHistOK}{yes}
\newcommand{\vgRealNPeriod}{20.0}
\newcommand{\vgRealNSigma}{1.94042}
\newcommand{\vgRealNLo}{15.84267}
\newcommand{\vgRealNHi}{20.0}
\newcommand{\vgRealNRayPeriod}{15.95265}
\newcommand{\vgRealNSpan}{504.0}
\newcommand{\vgRealNUnc}{1.94042}
\newcommand{\vgRealNSpread}{0.70181}
\newcommand{\vgRealNRecovers}{no}
\newcommand{\vgRealNHist}{16.11}
\newcommand{\vgRealNHistOK}{no}
\newcommand{\vgRealULamy}{17.247864}
\newcommand{\vgRealULamyOK}{yes}

\newcommand{\code}[1]{\texttt{#1}}
\shorttitle{A blind reanalysis of the ice-giant radio rotation periods}
\shortauthors{Barbere}
\begin{document}
\title{Non-Recovery of the Uranus and Neptune Rotation Periods in a Blind Reanalysis of
the Voyager~2 PRA Data}
\author[0009-0008-3289-4447]{Joseph Barbere}
\affiliation{Independent researcher}

\begin{abstract}
The rotation periods anchoring the System~III longitude systems of Uranus ($17.24\pm0.01$~h) and
Neptune ($16.11$~h) were fixed by the Voyager~2 Planetary Radio Astronomy (PRA) experiment during
the 1986 and 1989 flybys \citep{warwick1986,warwick1989}; the radio data have not since been
reanalysed blind with modern statistics, even as \citet{lamy2025} recently shifted Uranus's period
by 28~s from Hubble UV aurorae. We search the open PDS-PPI PRA low-band 6-second encounter volumes
\citep{pdsppi_vg2pra} with a blind Lomb--Scargle periodogram over a wide, un-tuned 14--20~h window,
using none of the original beaming or magnetic-longitude modelling, and take a rotation-block
bootstrap and the sub-band peak spread as the uncertainties. The result is a controlled null.
The pipeline is validated end-to-end: it recovers an injected
period in a synthetic flyby searched over the same window (\vgSynInjected~h $\to$
\vgSynRecovered~h) to $\sim$1~min but recovers neither real period: the Uranus peak is unstable across
sub-bands (\vgRealUPeriod~h, sub-band standard deviation \vgRealUSpread~h, offset $+1.2$~h from $17.24$~h; not
recovered, \vgRealURecovers), and the Neptune peak rails to the search bound (\vgRealNPeriod~h; not
recovered, \vgRealNRecovers). The method recovers a clean signal when one is present, so the
real-data failure is a limitation of blind total-power analysis, not of the pipeline: the
ice-giant total-power radio flux is not a clean rotational sinusoid, consistent with the geometric
modelling used in the original analyses having been necessary. This blind, total-power analysis places no useful bound on the period: its window resolution is
$\sim$\vgRealURes--\vgRealNRes~h and a peak unconstrained within the window; that reflects this
method on this data, not the flybys, since the modelled 1993 reanalysis reached $\pm 0.006$~h
($\pm 22$~s) from the same Neptune data, hundreds of times coarser than the Hubble result. The
contribution is a reproducible, modern-statistical measurement of how little a naive
total-power periodogram of these short flybys can constrain: injections into the real series
exclude a sinusoidal modulation at the historical periods only at depths of \vgRealUMinDepth\%
(Uranus) and \vgRealNMinDepth\% (Neptune).
\end{abstract}

\keywords{Planetary science (1255) --- Uranus (1751) --- Neptune (1096) --- Radio astronomy (1338)
--- Planetary magnetospheres (997) --- Time series analysis (1916)}

\section{Introduction}
The rotation period of a giant planet's deep interior is inferred from the periodic modulation of
magnetospheric phenomena locked to the internal magnetic field. For Uranus and Neptune the defining
measurement was the periodicity of auroral radio emission (Uranus Kilometric Radiation, UKR, and
its Neptune analogue), recorded by Voyager~2 PRA during the only spacecraft encounters with the ice
giants \citep{warwick1986,warwick1989}. Those periods, $17.24\pm0.01$~h and $16.11$~h, became the
System~III longitude systems still in use.

\citet{lamy2025} used 11~yr of Hubble UV aurora tracking to refine Uranus's period to
$\vgRealULamy\pm0.000010$~h, a 28~s shift and a $10^3$-fold precision gain over the Voyager value.
The radio value was a single brief-flyby measurement, revisited only by \citet{lecacheux1993}
(whose $16.108\pm0.006$~h refinement is itself finer than the modern Hubble shift) but never with a
blind, un-modelled search. In addition, the PRA refurbishing effort \citep{cecconi2017} reprocessed
only Jupiter and Saturn, leaving the ice-giant volumes an unworked archive; the original period
determinations were not simple periodograms but folded in magnetic-field models and emission
beaming. We therefore ask the complementary question: how much does a blind modern periodogram of
the raw radio flux, using none of that modelling, recover on its own? We answer it, with a
validated pipeline, in the negative.

\section{Data}
We use the open PDS-PPI PRA low-band 6-second RDR encounter volumes \citep{pdsppi_vg2pra}, datasets
\code{VG2-U-PRA-3-RDR-LOWBAND-6SEC-V1.0} (Uranus) and \code{VG2-N-PRA-3-RDR-LOWBAND-6SEC-V1.0}
(Neptune). Each is a fixed-width ASCII table of 48-second major frames; a frame carries a
spacecraft-event date and second-of-day plus eight sweeps of one status word and 70 channel
intensities in millibell, the sweeps spaced 6~s apart. The 70 low-band channels run from 1326.0 to
1.2~kHz in 19.2~kHz steps. Expanding the sweeps gives a 6-second-cadence dynamic spectrum spanning
$\sim$\vgRealUSpan~h at Uranus and $\sim$\vgRealNSpan~h at Neptune, of order 18--31 rotations.

\section{Method}
\label{sec:method}
For the flux series, we sum the linear power ($10^{\mathrm{mB}/1000}$ per channel) over the
100--1000~kHz band where the auroral emission dominates, giving one total-power sample per 6-second
sweep.

A rotation would modulate a continuous, red-noise-dominated emission, so
we search the (log, 0.1-h-binned, quadratically-detrended) flux with a blind Lomb--Scargle
periodogram over
a wide, physically-motivated 14--20~h ice-giant window, deliberately not tuned to bracket
the known values, since a narrow window centred on the answer would manufacture a match. Because the
samples are strongly autocorrelated and enormous in number, the analytic Lomb--Scargle false-alarm
probability is meaningless (absurdly small) and we do not report it. The statistical
uncertainty is a rotation-block bootstrap (resample whole rotation-length blocks); the systematic
one is the peak's spread across three disjoint thirds of the 100--1000~kHz analysis band. A coherent signal, whether rotation or a
viewing-geometry, beaming, or spectral-window modulation, is achromatic, so a small sub-band spread
is evidence of a real coherent peak but does not establish that it is the rotation. A discrete-epoch
Rayleigh $Z^2_1$ fold \citep{buccheri1983} of detected bursts, reused from this repository's FRB
module, is reported as a weaker cross-check.

We call the historical period \emph{recovered} only if the blind peak lands within one
combined-$\sigma$ of it, is band-robust (sub-band spread $<0.6$~h), and is not railed to a
search-window bound; proximity alone is worthless when $\sigma\sim2$~h.

\section{Synthetic recover-a-known}
We validate the pipeline on a synthetic flyby: a rotation-modulated flux series with an injected
period (\vgSynInjected~h), noise, and data gaps, searched in the same wide 14--20~h window
used on the real data. The Lomb--Scargle estimator recovers \vgSynRecovered~h (\vgSynRecoveredOK) and
the Rayleigh fold recovers \vgSynRayPeriod~h (\vgSynRayOK), both inside their bootstrap bands. This is
the control that makes the real-data null interpretable: the method finds a clean rotation when one is
present, so a failure on the real data reflects this method on this data rather than a broken
pipeline.

\section{Results}
\label{sec:results}
On the real encounter volumes the blind periodogram recovers neither rotation period
(Table~\ref{tab:periods}).

For Uranus, the peak sits at \vgRealUPeriod~h, offset $+1.2$~h from the historical
\vgRealUHist~h, or 1.1 window-resolution elements (\vgRealURes~h at this baseline), so
the offset alone is marginal. The verdict rests on achromaticity: across three disjoint
thirds of the 100--1000~kHz analysis band the peaks have a standard deviation of \vgRealUSpread~h (one railed),
so it is not a stable coherent line and not a recovery (\vgRealURecovers). (An earlier
version's band set included one band extending above the analysis band, which alone decided
this verdict; the disjoint in-band set reaches the same conclusion on independent channels.)
The peak is nonetheless significant against a scramble null
(FAP $\le \vgRealUFap$ over \vgRealUNIndep\ independent frequencies). A scramble also
destroys red noise and the flyby envelope, so this shows the series is not white noise, not
that it holds a coherent line; whatever structure there is, it is not at the rotation period. The historical and Hubble values
fall inside the $\pm\vgRealUUnc$~h band only because that band is nearly as wide as the
search window ($\sigma/W = 0.32$, against 0.29 for a uniform distribution); the bootstrap
is reporting an unconstrained peak, not an uncertainty, and we do not quote it as one.

For Neptune, the peak rails to the upper search bound (\vgRealNPeriod~h), a bound,
not an estimate, so no consistency statement with the historical $\vgRealNHist \pm 0.05$~h is
possible (the macro reports ``\vgRealNHistOK''). The offset from 16.11~h spans several resolution
elements (\vgRealNRes~h each at this baseline), a real mismatch; the railing is the signature of a periodogram dominated by the
red-noise / flyby-envelope continuum rather than a rotational line. We disclose the analysis
history in full, since this paper's title contains ``blind'': an earlier version of this
analysis used a narrower search window and reported Neptune as recovered near 16~h; widening
the window moved the peak to the bound, exposing the earlier result as a window artifact, and
the recovery criterion's band-robustness and railing clauses were added after that episode. A
second change followed: the sub-band set was redefined as disjoint thirds after the original set
was found to decide the Uranus verdict (above). The scramble and injection controls below were
also added after the real result and are post hoc. The final window and criterion were applied
identically to both planets and to the control. A narrower window still places the Neptune peak near 16~h; that
demonstrates window-dependence, not a rotation.

We tested the selection function directly. Injecting a fractional sinusoid at each
historical period into the real flux series (red noise, flyby envelope and detrend included)
and re-running the blind search recovers it in every one of eight phases at depths of
\vgRealUMinDepth\% (Uranus) and \vgRealNMinDepth\% (Neptune), the smallest tested depths to reach
that on a grid of 2, 5, 10, 20 and 40\%, so the true eight-of-eight threshold lies between half
and all of each value. A sinusoidal total-power modulation at the historical period is excluded at
those depths. This, not the
synthetic spike train, is the operative validation: the synthetic control (which passes the
full three-clause criterion, \vgSynFullOK, with seed-to-seed scatter \vgSynSeedSd~h across
\vgSynSeedN\ seeds) exercises the machinery end to end, while the injection measures what this search could have seen.

The blind total-power periodogram does not recover the ice-giant rotation
periods, while recovering a clean injected signal in the identical analysis. The auroral total-power
flux is therefore not a clean rotational sinusoid over these flybys, consistent with the historical
determinations' beaming and magnetic-longitude modelling, which a blind periodogram omits, having been
necessary; this analysis does not test that modelling. The secondary Rayleigh burst fold (Uranus \vgRealURayPeriod~h,
Neptune \vgRealNRayPeriod~h) is likewise uninformative, as expected for a discrete-epoch method on
this emission.

\begin{deluxetable}{lccc}
\tablecaption{Rotation periods (hours)\label{tab:periods}}
\tablehead{\colhead{Source} & \colhead{Uranus} & \colhead{Neptune} & \colhead{method}}
\startdata
Voyager (historical)   & $17.24\pm0.01$ & $16.11$ & radio (modelled) \\
\citet{lamy2025}       & $\vgRealULamy$ & \nodata & HST UV \\
This work (blind)      & $\vgRealUPeriod$ & $\vgRealNPeriod$ (railed) & radio (blind LS) \\
recovers historical?   & \vgRealURecovers & \vgRealNRecovers & \\
\enddata
\tablecomments{This work: blind Lomb--Scargle of the PRA flux series over an un-tuned 14--20~h window;
uncertainty = larger of the rotation-block bootstrap $\sigma$ and the sub-band spread. ``Recovered''
requires a band-robust, un-railed peak within one combined-$\sigma$ of the historical value. Neither
planet qualifies.}
\end{deluxetable}

\section{Reproducibility}
The metric, its synthetic recover-a-known, and all unit tests (including a parser test against a
vendored real PRA sample) run offline in CI (\code{python -m jansky\_research.vgpra --offline}). The
real reanalysis downloads the two open PDS-PPI encounter volumes and runs on demand
(\code{scripts/vgpra\_real.py}). Code, tests, and this paper's pipeline are in the
\code{jansky-research} repository, built on the \code{jansky} course library \citep{jansky} and its
\code{frbperiod} periodicity machinery. Draft prepared with the assistance of a large language model
\citep{anthropicclaude}; all numerical results are pipeline-generated and human-verified.

\software{\code{jansky-research} \citep{janskyresearch}, \code{jansky} \citep{jansky},
NumPy, SciPy, Astropy. Analysis assisted by Claude \citep{anthropicclaude}; all results
verified by the pipeline's test suite.}

\bibliographystyle{aasjournal}
\bibliography{refs}
\end{document}
